\documentclass[11pt,a4paper]{article}

\usepackage[margin=1in]{geometry}
\usepackage{amsmath,amssymb,bm}
\usepackage{booktabs}
\usepackage{graphicx}
\usepackage{hyperref}
\usepackage{xcolor}
\usepackage{fancyhdr}
\usepackage{enumitem}
\usepackage{siunitx}
\usepackage{float}
\usepackage{titlesec}
\usepackage{parskip}

\hypersetup{colorlinks=true, linkcolor=blue!70!black, urlcolor=blue!70!black, citecolor=blue!70!black}

\pagestyle{fancy}
\fancyhead[L]{\small\textit{SEAMOTH Glide Theory}}
\fancyhead[R]{\small\textit{Middle-Aged Engineering}}
\fancyfoot[C]{\thepage}

\newcommand{\seamoth}{\textsc{Seamoth}}
\newcommand{\half}{\tfrac{1}{2}}
\newcommand{\vect}[1]{\bm{#1}}
\newcommand{\mat}[1]{\mathbf{#1}}
\newcommand{\deriv}[2]{\frac{\mathrm{d}#1}{\mathrm{d}#2}}
\newcommand{\pderiv}[2]{\frac{\partial #1}{\partial #2}}

\newcommand{\note}[1]{\par\smallskip\noindent\colorbox{blue!8}{\parbox{\dimexpr\linewidth-2\fboxsep}{%
  \small\textbf{SEAMOTH:} #1}}\smallskip}

\newcommand{\insight}[1]{\par\smallskip\noindent\colorbox{green!8}{\parbox{\dimexpr\linewidth-2\fboxsep}{%
  \small\textbf{Physical insight:} #1}}\smallskip}

\newcommand{\warning}[1]{\par\smallskip\noindent\colorbox{red!8}{\parbox{\dimexpr\linewidth-2\fboxsep}{%
  \small\textbf{Warning:} #1}}\smallskip}

\title{\vspace{-1cm}\textbf{Underwater Glider Theory\\for the \seamoth{} Project}\\[8pt]
\large A comprehensive technical reference following Graver (2005),\\
applied to a shallow-water vectored-buoyancy-control glider}
\author{Middle-Aged Engineering}
\date{March 2026 \quad---\quad v0.1}

\begin{document}
\maketitle

\begin{abstract}
This document extracts and explains every piece of theory from Graver's
\emph{Underwater Gliders: Dynamics, Control and Design} (Princeton, 2005) that is
relevant to the \seamoth{} v1 reef-survey glider.  It is intended to be read
as a self-contained tutorial: we derive the equations from first principles, explain
the physical intuition behind each result, and plug in \seamoth{} parameters throughout
so the reader can cross-check every number against the companion Jupyter notebook.
Where Graver's framework does not cover \seamoth{}'s unique features (notably the
4-bladder vectored buoyancy control), we note the gap and describe how to extend
the theory.

This is \emph{not} a summary of the dissertation.  It is a guided tour through the
parts that matter for designing, building, and flying our glider, written so that
a team member with undergraduate-level mechanics and fluid dynamics can follow every step.

Cross-references to Graver are given as (G-\S$n$) for section $n$ and (G-Eq.$\,m$) for
equation $m$.
\end{abstract}

\tableofcontents
\newpage

%======================================================================
%  PART I: FOUNDATIONS
%======================================================================
\part{Foundations}

%----------------------------------------------------------------------
\section{What is an underwater glider?}
\label{sec:intro}
%----------------------------------------------------------------------

An underwater glider is an autonomous vehicle that moves through the ocean by
controlling its buoyancy.  It has no propeller.  Wings convert the resulting vertical
motion (sinking or rising) into forward glide, exactly like a sailplane---except
the vehicle is designed to be neutrally buoyant in a fluid 800 times denser than air.

The flight cycle is a sawtooth: the vehicle makes itself heavy (negative buoyancy),
glides downward and forward, then makes itself light (positive buoyancy) and glides
upward and forward.  At the surface it gets a GPS fix and communicates; then it dives
again.

\seamoth{} differs from conventional gliders (Slocum, Spray, Seaglider) in three ways:
\begin{enumerate}[nosep]
  \item \textbf{Vectored buoyancy control (VBC):} four independent bladders instead
    of one, providing pitch \emph{and} roll authority without a sliding internal mass.
  \item \textbf{Shallow water:} 5--30\,m operating depth, versus 200--1500\,m for
    ocean gliders.  This means higher cycling frequency and different energy trade-offs.
  \item \textbf{Agility:} target turning radius $<$5\,m and pitch transitions $<$5\,s,
    versus minutes for ocean gliders.
\end{enumerate}

The theory in this document applies to all buoyancy-driven gliders.  We note
where \seamoth{}'s VBC architecture requires extensions beyond Graver's single-bladder model.


%----------------------------------------------------------------------
\section{Coordinate frames and kinematics}
\label{sec:kinematics}
%----------------------------------------------------------------------

Three reference frames are needed to describe glider motion (G-\S3.1).

\subsection{Inertial frame}

An earth-fixed frame with orthogonal axes $(x,y,z)$ and unit vectors
$(\vect{i},\vect{j},\vect{k})$.  The $x$-$y$ plane is horizontal; the $z$-axis
points \emph{downward} (positive $z$ = increasing depth).  This is the standard
convention in marine vehicle dynamics.

\warning{The $z$-down convention means that ``going up'' is \emph{negative}~$\dot{z}$.
This catches people who are used to aerospace $z$-up frames.  Our IEKF and MPC firmware
must use the same convention consistently.}

\subsection{Body frame}

A frame fixed to the glider with its origin at the \textbf{centre of buoyancy} (CB)
and axes aligned with the vehicle's principal axes:
\begin{itemize}[nosep]
  \item $e_1$: forward (nose direction),
  \item $e_2$: to the right (starboard wing),
  \item $e_3$: downward (through the belly).
\end{itemize}

The orientation of the body frame relative to the inertial frame is described by the
rotation matrix $\mat{R}$, parameterized by three Euler angles in the standard
yaw--pitch--roll (ZYX) convention:
\begin{align}
  \psi &\quad\text{yaw (heading), positive clockwise from above}, \\
  \theta &\quad\text{pitch, positive nose-up}, \\
  \phi &\quad\text{roll, positive right-wing-down}.
\end{align}

The kinematic equations are (G-Eq.\,3.1--3.2):
\begin{align}
  \dot{\mat{R}} &= \mat{R}\,\hat{\vect{\Omega}}, \label{eq:Rdot} \\
  \dot{\vect{b}} &= \mat{R}\,\vect{v}, \label{eq:bdot}
\end{align}
where $\vect{b}=(x,y,z)^T$ is the glider position, $\vect{v}=(v_1,v_2,v_3)^T$ is the
translational velocity in the body frame, $\vect{\Omega}=(\Omega_1,\Omega_2,\Omega_3)^T$
is the angular velocity in the body frame, and $\hat{\cdot}$ denotes the skew-symmetric
matrix representing a cross product:
\begin{equation}
  \hat{\vect{x}} = \begin{pmatrix} 0 & -x_3 & x_2 \\ x_3 & 0 & -x_1 \\ -x_2 & x_1 & 0 \end{pmatrix}
  \quad\text{so that}\quad \hat{\vect{x}}\,\vect{y} = \vect{x}\times\vect{y}.
\end{equation}

\subsection{Wind frame}

The wind frame tracks the glider's velocity \emph{relative to the water},
$\vect{v}_r$.  In the absence of currents, $\vect{v}_r = \vect{v}$.
Its orientation relative to the body frame is described by two
\textbf{aerodynamic angles} (G-\S3.1.3):

\begin{align}
  \alpha &= \arctan\!\left(\frac{v_{r3}}{v_{r1}}\right)
    &&\text{angle of attack}, \label{eq:aoa} \\
  \beta  &= \arctan\!\left(\frac{v_{r2}}{|\vect{v}_r|}\right)
    &&\text{sideslip angle}. \label{eq:sideslip}
\end{align}

\insight{The angle of attack $\alpha$ is the angle between the glider's nose ($e_1$)
and its actual direction of travel through the water.  It is the wing's ``working angle''---the
parameter that determines how much lift the wing generates.  Sideslip $\beta$ is the
lateral equivalent: nonzero $\beta$ means the glider is ``crabbing'' sideways through the water.}

\subsection{The glide-path angle $\xi$}
\label{sec:xi}

The \textbf{glide-path angle} is the angle of the velocity vector relative to the
horizontal:
\begin{equation}
  \xi = \theta - \alpha.
  \label{eq:xi}
\end{equation}

Graver's sign convention: $\xi < 0$ for a downward glide, $\xi > 0$ for an upward glide.

\warning{Throughout this document and in the Jupyter notebook, we quote $\xi$ as a
positive number in degrees (e.g.\ ``$\xi=15°$'' meaning $15°$ below horizontal).
The equilibrium solver converts to Graver's convention internally.  Always check which
convention a particular function expects.}

\subsection{Currents}

In the presence of a current $\vect{V}_c$ (expressed in the inertial frame), the
velocity relative to the water is:
\begin{equation}
  \vect{v}_r = \vect{v} - \mat{R}^T \vect{V}_c.
  \label{eq:vr}
\end{equation}

All hydrodynamic forces depend on $\vect{v}_r$, not $\vect{v}$.  For equilibrium
analysis we assume still water ($\vect{V}_c=\vect{0}$), but in the Arabian Gulf
tidal currents of 0.1--0.3\,m/s are significant relative to our 0.5\,m/s glide speed.
The IEKF must estimate current as a slowly-varying state.


%----------------------------------------------------------------------
\section{The vehicle model}
\label{sec:vehicle}
%----------------------------------------------------------------------

\subsection{Mass decomposition}

Graver decomposes the glider mass into several components (G-\S3.2.1):

\begin{description}[nosep]
  \item[$m_h$:] Hull mass (uniformly distributed structural mass).
  \item[$\bar{m}$:] Movable point mass (the sliding battery pack that controls
    pitch and roll in conventional gliders).
  \item[$m_b$:] Variable ballast mass (the oil pumped in/out to change buoyancy),
    located at CB.
  \item[$m_w$:] Static offset mass (accounts for non-uniform weight distribution).
  \item[$m_s$:] Stationary mass: $m_s = m_h + m_w + m_b$.
  \item[$m_v$:] Total vehicle mass: $m_v = m_s + \bar{m}$.
  \item[$m$:] Mass of displaced water: $m = \rho\,V_d$ where $V_d$ is the
    displaced volume.
  \item[$m_0$:] \textbf{Excess mass}: $m_0 = m_v - m$.
\end{description}

The excess mass $m_0$ is the key control variable.  When $m_0 > 0$ the glider is
heavier than the water it displaces and sinks; when $m_0 < 0$ it rises.  The buoyancy
engine controls $m_0$ by moving oil between an internal reservoir and external bladders,
changing the vehicle's displaced volume $V_d$ without changing its actual mass $m_v$.

\insight{This is a subtle but important point.  The pump doesn't add or remove mass
from the vehicle---it changes the vehicle's \emph{volume}.  Pumping oil into the external
bladders makes the vehicle effectively larger (more water displaced $\rightarrow$ $m$
increases $\rightarrow$ $m_0$ decreases $\rightarrow$ more buoyant).  Retracting oil
makes the vehicle smaller and heavier.}

\note{\seamoth{} parameters: $m_v\approx8.5$\,kg, $V_d\approx8.3$\,L,
$m=\rho V_d=1023\times0.0083\approx8.49$\,kg.  At neutral trim $m_0\approx0$.
The buoyancy engine can vary $V_d$ by $\pm160$\,mL, giving a maximum
$|m_0|\approx\rho\times160\times10^{-6}\approx0.164$\,kg.  In the notebook, $m_0$
is input directly in kg.}

\subsection{Application to \seamoth{} VBC}

\seamoth{} does \emph{not} have a sliding mass $\bar{m}$.  Instead, four independent
bladders provide both buoyancy control and attitude control:

\begin{itemize}[nosep]
  \item \textbf{Pitch:} inflate forward bladders while deflating aft bladders (or vice
    versa).  This shifts the centre of buoyancy (CB) forward or aft of the centre of
    gravity (CG), creating a pitch moment.
  \item \textbf{Roll:} inflate port bladders while deflating starboard bladders.
    This shifts CB laterally, creating a roll moment.
  \item \textbf{Net buoyancy:} inflate or deflate all bladders together to change $m_0$.
\end{itemize}

Graver's model uses a single $m_0$ for net buoyancy and a sliding mass position
$\vect{r}_p$ for attitude control.  For \seamoth{}, we can map the 4-bladder state
into equivalent $m_0$ (total) and equivalent CB shift (attitude), but the detailed
dynamics of VBC are beyond Graver's framework.  The ReefGlider paper (Galloway et al.,
2024) provides the VBC-specific theory.

For the purpose of \emph{steady-state} analysis (which is what this document covers),
the VBC vs.\ sliding-mass distinction does not matter: at equilibrium, all that matters
is the total $m_0$ and the resulting pitch angle $\theta$.


%----------------------------------------------------------------------
\section{Forces: gravity, buoyancy, and hydrodynamics}
\label{sec:forces}
%----------------------------------------------------------------------

\subsection{Gravity and buoyancy}

The gravitational force on the vehicle is $m_v\,g$ downward.  The buoyancy force is
$m\,g = \rho\,V_d\,g$ upward.  The \emph{net} gravitational--buoyancy force is:
\begin{equation}
  F_\text{net} = m_0\,g = (m_v - m)\,g,
  \label{eq:Fnet}
\end{equation}
acting downward when $m_0>0$ (negative buoyancy).

In the body frame, this force has components that depend on the pitch angle $\theta$:
\begin{equation}
  \vect{F}_\text{grav,body} = m_0\,g\begin{pmatrix} -\sin\theta \\ 0 \\ \cos\theta \end{pmatrix}.
  \label{eq:Fgrav_body}
\end{equation}

\insight{When the glider pitches nose-down ($\theta<0$ for a downward glide), the
$e_1$ component of gravity is \emph{positive} (forward)---this is what drives the glide.
The $e_3$ component pushes the glider ``down'' in the body frame, which is resisted by lift.}

\subsection{Hydrodynamic forces}

The fluid exerts forces and moments on the glider.  These are decomposed into
(G-\S3.2.7):

\begin{description}[nosep]
  \item[Lift $L$:] Force perpendicular to $\vect{v}_r$ (in the $e_1$-$e_3$ plane),
    directed from the belly toward the back of the glider.
  \item[Drag $D$:] Force opposing $\vect{v}_r$.
  \item[Side force $SF$:] Force perpendicular to $\vect{v}_r$ in the lateral direction
    (zero in symmetric flight).
  \item[Moments:] Pitch moment $M$, roll moment $\mathcal{L}$, yaw moment $N$.
\end{description}

In addition to these ``viscous'' forces, the dense fluid creates \textbf{added mass}
effects: accelerating through water requires accelerating the surrounding fluid as well.
The added mass matrix $\mat{M}_f$ and added inertia matrix $\mat{J}_f$ appear in the
full equations of motion (G-\S3.2.3).  At steady state (constant velocity and angular
velocity), added mass terms vanish---they only matter during transients like inflections.

\subsection{Quasi-steady hydrodynamic model}

For steady-state analysis, Graver uses a coefficient-based model where lift, drag,
and moment are functions of the instantaneous angle of attack $\alpha$ and speed $V$:
\begin{align}
  L &= \half\rho\,V_r^2\,S\,C_L(\alpha), \label{eq:L} \\
  D &= \half\rho\,V_r^2\,S\,C_D(\alpha), \label{eq:D} \\
  M &= \half\rho\,V_r^2\,S\,c\,C_M(\alpha), \label{eq:M}
\end{align}
where $S$ is the wing reference area, $c$ is the wing chord, and $V_r = |\vect{v}_r|$.

The coefficient models are:
\begin{align}
  C_L &= C_{L0} + C_{L\alpha}\,\alpha, \label{eq:CL} \\
  C_D &= C_{D0} + C_{D\alpha}\,\alpha^2, \label{eq:CD} \\
  C_M &= C_{M0} + C_{M\alpha}\,\alpha. \label{eq:CM}
\end{align}

These are the simplest polynomial fits that capture the essential physics: lift increases
linearly with $\alpha$ (until stall), drag has a quadratic ``drag bucket'' centred near
$\alpha=0$, and the pitch moment is approximately linear in $\alpha$.


%======================================================================
%  PART II: BUILDING THE HYDRODYNAMIC MODEL
%======================================================================
\part{Building the Hydrodynamic Model for \seamoth{}}

%----------------------------------------------------------------------
\section{Lift model}
\label{sec:lift}
%----------------------------------------------------------------------

\subsection{2D airfoil data}

The SD7003 airfoil (Selig, UIUC) was designed for low Reynolds number applications
and has extensive wind-tunnel data at $\mathrm{Re}=40{,}000$--$200{,}000$.
At our operating Reynolds numbers ($\sim$40k--60k), the key parameters are:

\begin{itemize}[nosep]
  \item 2D lift-curve slope: $C_{l\alpha,\text{2D}} \approx 6.2\;\mathrm{rad}^{-1}$
    (close to thin-airfoil theory's $2\pi\approx6.28$, slightly reduced by viscous effects).
  \item Zero-$\alpha$ lift: $C_{l0,\text{2D}} \approx 0.15$ (due to the airfoil's camber).
  \item Maximum $C_l$ before stall: $\sim$0.9 at $\alpha\approx8$--$9°$ (Re\,50k).
\end{itemize}

\subsection{Finite wing correction}

A real wing of finite span generates tip vortices that reduce the effective angle of
attack and the lift-curve slope.  The Prandtl lifting-line correction
gives (G-\S3.2.7):
\begin{equation}
  C_{L\alpha,\text{3D}} = C_{l\alpha,\text{2D}}\;\frac{\mathrm{AR}}{\mathrm{AR}+2},
  \label{eq:CLa3d}
\end{equation}
where $\mathrm{AR}=b^2/S$ is the aspect ratio ($b$ = wingspan).

\insight{The factor $\mathrm{AR}/(\mathrm{AR}+2)$ is an approximation of the exact
Prandtl result, valid for wings with elliptic or near-elliptic loading.  For
$\mathrm{AR}=6$, it gives 75\% of the 2D value.  Higher AR $\rightarrow$ less tip
loss $\rightarrow$ more lift per radian of $\alpha$.  This is why sailplanes have
long, slender wings.}

\subsection{Hull lift contribution}

The glider hull (a slender body of revolution) contributes additional lift through
a combination of crossflow force and interference effects with the wing.  Graver
accounts for this with a multiplicative factor (G-\S5.2.1):
\begin{equation}
  C_{L\alpha} = C_{L\alpha,\text{3D}} \times k_h,
  \label{eq:CLa_total}
\end{equation}
where $k_h \approx 1.05$--$1.15$ depending on hull fineness ratio and wing position.

\note{$C_{l\alpha,\text{2D}}=6.2$, $\mathrm{AR}=6$, $k_h=1.10$:
\begin{align*}
  C_{L\alpha,\text{3D}} &= 6.2 \times \frac{6}{6+2} = 6.2 \times 0.75 = 4.65\;\mathrm{rad}^{-1}, \\
  C_{L\alpha} &= 4.65 \times 1.10 = \mathbf{5.115}\;\mathrm{rad}^{-1}.
\end{align*}
$C_{L0} = 0.15$ (assumed unchanged by finite-wing effects, as camber lift is primarily 2D).}


%----------------------------------------------------------------------
\section{Drag model}
\label{sec:drag}
%----------------------------------------------------------------------

The total drag coefficient (referenced to wing area $S$) is the sum of several
components:
\begin{equation}
  C_D = \underbrace{C_{D0}}_{\text{parasitic}} + \underbrace{C_{D\alpha}\,\alpha^2}_{\text{induced}}.
  \label{eq:CDtotal}
\end{equation}

\subsection{Parasitic drag $C_{D0}$}

This is all drag at zero lift (or more precisely, at $\alpha=0$).  It includes:

\paragraph{Hull drag.}
The hull is a bluff body.  Its drag coefficient is typically referenced to
\emph{frontal area} $A_f$:
\begin{equation}
  D_\text{hull} = \half\rho\,V^2\,A_f\,C_{d,\text{hull}}.
\end{equation}
To reference this to wing area $S$, multiply by $A_f/S$:
\begin{equation}
  C_{D0,\text{hull}} = C_{d,\text{hull}}\;\frac{A_f}{S}.
  \label{eq:cd0hull}
\end{equation}

The hull $C_d$ depends strongly on the nose shape, tail fairing, and surface finish.
A bare cylinder has $C_d\approx0.8$--$1.0$; a well-faired torpedo shape can achieve
$C_d\approx0.04$--$0.08$.  Our 4'' BlueRobotics tube with a moulded nose cone and
tail fairing is estimated at $C_d\approx0.11$.

\warning{The hull frontal area must use the \emph{outer} diameter (114.3\,mm), not the
inner diameter ($\sim$100\,mm).  The O-ring flanges and end caps may add additional
frontal area.  This was the source of the error in our original spec document.}

\note{$A_f = \pi/4 \times 0.1143^2 = 102.6\;\mathrm{cm}^2 = 0.01026\;\mathrm{m}^2$.\\
$C_{D0,\text{hull}} = 0.11 \times (0.01026/0.030) = 0.11 \times 0.342 = \mathbf{0.0376}$.}

\paragraph{Wing profile drag.}
The wing's own friction and pressure drag at zero lift.  From SD7003 wind tunnel data
at Re\,50k: $C_{d0,\text{wing}} \approx 0.025$.

\paragraph{Appendage drag.}
Tail surfaces, antenna mast, fairings, cable penetrators, and any other protuberances.
Estimated as $C_{d,\text{app}} \approx 0.008$ (referenced to $S$).

\paragraph{Total parasitic drag.}
\begin{equation}
  C_{D0} = C_{D0,\text{hull}} + C_{d0,\text{wing}} + C_{d,\text{app}}.
  \label{eq:CD0total}
\end{equation}

\note{$C_{D0} = 0.0376 + 0.025 + 0.008 = \mathbf{0.0706}$.\\
The hull contributes 53\% of total parasitic drag---this is the single biggest drag source
on the vehicle and the primary target for aerodynamic improvement.}

\subsection{Induced drag $C_{D\alpha}$}

When the wing generates lift, it also generates lift-induced drag through the tip
vortex system.  From lifting-line theory with Oswald span efficiency factor $e$:
\begin{equation}
  C_{Di} = \frac{C_L^2}{\pi\,\mathrm{AR}\,e}.
  \label{eq:CDi}
\end{equation}

For the drag polar $C_D = C_{D0} + C_{D\alpha}\alpha^2$, we substitute
$C_L \approx C_{L\alpha}\alpha$ (ignoring $C_{L0}$ for the $\alpha$-dependent term):
\begin{equation}
  C_{D\alpha} = \frac{C_{L\alpha}^2}{\pi\,\mathrm{AR}\,e}.
  \label{eq:CDalpha}
\end{equation}

\note{$C_{D\alpha} = 5.115^2 / (\pi \times 6 \times 0.85) = 26.16/16.02 = \mathbf{1.633}\;\mathrm{rad}^{-2}$.}

\subsection{Drag breakdown at operating conditions}

At a typical operating point ($V=0.53$\,m/s, $\alpha=1.3°=0.023$\,rad):
\begin{align*}
  q &= \half\times1023\times0.53^2 = 143.7\;\text{Pa}, \\
  D_\text{hull} &= 143.7\times0.01026\times0.11 = 0.162\;\text{N} \quad(53\%), \\
  D_\text{wing,profile} &= 143.7\times0.030\times0.025 = 0.108\;\text{N} \quad(35\%), \\
  D_\text{wing,induced} &= 143.7\times0.030\times1.633\times0.023^2 = 0.004\;\text{N} \quad(1\%), \\
  D_\text{appendages} &= 143.7\times0.030\times0.008 = 0.035\;\text{N} \quad(11\%), \\
  D_\text{total} &= 0.308\;\text{N}.
\end{align*}

\insight{At our low angles of attack, induced drag is negligible (only 1\% of total).
This is very different from sailplanes, where induced drag dominates.  The reason is
simple: at $\alpha=1.3°$, the wing barely generates any lift.  The hull's parasitic drag
overwhelms everything else.  This means that wing AR and Oswald efficiency are
\emph{not} critical design parameters for \seamoth{}---hull fairing is far more important.}


%----------------------------------------------------------------------
\section{Graver's lumped $K$ parameters}
\label{sec:Kparams}
%----------------------------------------------------------------------

To simplify the equilibrium equations, Graver defines four lumped parameters that
absorb the dynamic-pressure-times-area product (G-Eq.\,4.14--4.16):
\begin{align}
  K_{D0} &= \half\rho\,S\,C_{D0}, &
  K_D    &= \half\rho\,S\,C_{D\alpha}, \label{eq:KDdef} \\
  K_{L0} &= \half\rho\,S\,C_{L0}, &
  K_L    &= \half\rho\,S\,C_{L\alpha}. \label{eq:KLdef}
\end{align}

These have units of $\mathrm{N\cdot s^2/m^2}$ (force per velocity-squared).
The total hydrodynamic forces become:
\begin{align}
  D &= (K_{D0} + K_D\,\alpha^2)\,V^2, \label{eq:Dlumped} \\
  L &= (K_{L0} + K_L\,\alpha)\,V^2. \label{eq:Llumped}
\end{align}

\note{$\half\rho\,S = \half\times1023\times0.030 = 15.345\;\text{N}\cdot\text{s}^2/\text{m}^4$.
\begin{center}
\begin{tabular}{@{}ll@{}}
\toprule
$K_{D0}$ & $15.345\times0.0706 = \mathbf{1.083}$ \\
$K_D$    & $15.345\times1.633 = \mathbf{25.06}$ \\
$K_{L0}$ & $15.345\times0.150 = \mathbf{2.302}$ \\
$K_L$    & $15.345\times5.115 = \mathbf{78.49}$ \\
\bottomrule
\end{tabular}
\end{center}}


%----------------------------------------------------------------------
\section{Reynolds number regime}
\label{sec:reynolds}
%----------------------------------------------------------------------

The Reynolds number governs the character of the flow over the wing:
\begin{equation}
  \mathrm{Re}_c = \frac{V\,c}{\nu},
  \label{eq:Re}
\end{equation}
where $c$ is the chord length and $\nu$ is the kinematic viscosity.

\subsection{Temperature dependence of viscosity}

Seawater viscosity decreases with temperature.  This is significant because
\seamoth{} operates in warm Gulf water:

\begin{center}
\begin{tabular}{@{}lcc@{}}
\toprule
Environment & $T$ [\textdegree C] & $\nu$ [$\times10^{-6}\;\text{m}^2/\text{s}$] \\
\midrule
North Atlantic (Slocum) & 10--15 & 1.14--1.35 \\
Arabian Gulf (SEAMOTH) & 25--35 & 0.76--0.93 \\
Design point (28\textdegree C) & 28 & 0.88 \\
\bottomrule
\end{tabular}
\end{center}

The 30\% reduction in viscosity means that, at the same speed, \seamoth{}'s wings
operate at 30\% higher Reynolds number than a glider in cold water.  This pushes us
deeper into the SD7003's effective zone and away from the laminar-separation-bubble
(LSB) danger region below Re\,30k.

\note{At $V=0.53$\,m/s (design point, $\xi=15°$, $m_0=120$\,mL): $\mathrm{Re}_c = 0.53\times0.071/0.88\times10^{-6} = 42{,}800$.  Safely in the SD7003's optimal range.}

\subsection{Laminar separation bubbles}

Below $\mathrm{Re}\approx30{,}000$, the laminar boundary layer on the suction side
of the airfoil separates before transitioning to turbulence.  A ``bubble'' of
recirculating flow forms, dramatically increasing drag and often reducing lift.  This
sets a practical minimum speed:
\begin{equation}
  V_\text{min} = \frac{\mathrm{Re}_\text{min}\times\nu}{c}
  = \frac{30{,}000\times0.88\times10^{-6}}{0.071}
  = 0.37\;\text{m/s}.
\end{equation}

The buoyancy engine should always maintain enough ballast to exceed this speed.
At $\xi=15°$, the minimum ballast for $V>0.37$\,m/s is approximately 60\,mL.


%======================================================================
%  PART III: STEADY-STATE ANALYSIS
%======================================================================
\part{Steady-State Glide Analysis}

%----------------------------------------------------------------------
\section{The longitudinal (vertical plane) model}
\label{sec:longitudinal}
%----------------------------------------------------------------------

For most analysis, we restrict the glider to the vertical plane:
$\beta=0$ (no sideslip), $\phi=0$ (no roll), and all lateral velocities and
angular rates are zero.  This is the \textbf{longitudinal model} (G-\S4.1).

The remaining states are:
\begin{itemize}[nosep]
  \item Position: depth $z$ and horizontal distance $x$.
  \item Velocity: $V$ (speed through the water) and $\xi$ (glide-path angle).
  \item Attitude: pitch $\theta=\xi+\alpha$.
  \item Internal: movable mass position $r_{p1}$ (or VBC bladder state) and ballast $m_0$.
\end{itemize}

\subsection{Equations of motion in the vertical plane}

At steady state (constant $V$, $\xi$, $\alpha$), all accelerations vanish and the
equations of motion reduce to a force balance (G-\S4.1.2).  Resolving forces along
the velocity vector ($w_1$ axis) and perpendicular to it ($w_3$ axis):

\paragraph{Along the velocity:}
\begin{equation}
  -D - m_0\,g\,\sin\xi = 0
  \quad\Rightarrow\quad
  D = -m_0\,g\,\sin\xi.
  \label{eq:force_along}
\end{equation}

For a downward glide ($\xi<0$ in Graver's convention), $\sin\xi<0$ so the right side
is positive---gravity overcomes drag.

\paragraph{Perpendicular to the velocity:}
\begin{equation}
  L - m_0\,g\,\cos\xi = 0
  \quad\Rightarrow\quad
  L = m_0\,g\,\cos\xi.
  \label{eq:force_perp}
\end{equation}

Lift supports the ``sideways'' component of the net weight.

Dividing Eq.\,\eqref{eq:force_along} by Eq.\,\eqref{eq:force_perp}:
\begin{equation}
  \boxed{\frac{D}{L} = -\tan\xi}
  \label{eq:DoverL}
\end{equation}

This is the fundamental glide relation: the tangent of the glide angle equals the
drag-to-lift ratio.  A glider with $L/D=3.73$ glides at
$\xi=\arctan(1/3.73)=15.0°$.

\insight{This equation is identical for sailplanes and underwater gliders---it's
pure geometry.  The difference is that a sailplane's $L/D$ is typically 30--60
(shallow glide), while an underwater glider's $L/D$ is 2--6 (steep glide).  The
underwater glider's high drag comes from the bluff hull, which a sailplane doesn't have.}


%----------------------------------------------------------------------
\section{Equilibrium glide: solving for $\alpha$ and $V$}
\label{sec:equilibrium}
%----------------------------------------------------------------------

\subsection{Solving for $\alpha_d$}

Substituting the lumped force expressions (Eqs.\,\ref{eq:Dlumped}--\ref{eq:Llumped})
into the force balance:
\begin{align}
  (K_{D0} + K_D\alpha^2)\,V^2 &= -m_0\,g\sin\xi, \label{eq:fb1} \\
  (K_{L0} + K_L\alpha)\,V^2 &= m_0\,g\cos\xi. \label{eq:fb2}
\end{align}

Dividing to eliminate $V^2$:
\begin{equation}
  \frac{K_{D0} + K_D\alpha^2}{K_{L0} + K_L\alpha} = -\tan\xi.
  \label{eq:alphabalance}
\end{equation}

Rearranging into a quadratic in $\alpha$ (G-Eq.\,4.21):
\begin{equation}
  K_D\,\alpha^2 - K_L\tan\xi\;\alpha + (K_{D0} - K_{L0}\tan\xi) = 0.
  \label{eq:alphaquad}
\end{equation}

Applying the quadratic formula and selecting the physically meaningful root
(small, positive $\alpha$):
\begin{equation}
  \boxed{
    \alpha_d = \frac{K_L}{2\,K_D}\,\tan\xi\;
    \left(-1 + \sqrt{1 - \frac{4\,K_D}{K_L^2}\,\cot\xi\,(K_{D0}\cot\xi + K_{L0})}\right)
  }
  \label{eq:alpha_d}
\end{equation}

where $\xi$ carries Graver's sign convention (negative for downward glides).

\note{Numerical evaluation at $\xi=15°$ (Graver: $\xi=-15°$):
\begin{align*}
  \tan\xi &= -0.2679, \quad \cot\xi = -3.732 \\
  \text{disc.} &= 1 - \frac{4\times25.06}{78.49^2}\times(-3.732)\times(1.083\times(-3.732)+2.302) \\
  &= 1 - 0.01625\times(-3.732)\times(-1.739) \\
  &= 1 - 0.1054 = 0.8946 \\
  \alpha_d &= \frac{78.49}{2\times25.06}\times(-0.2679)\times(-1+\sqrt{0.8946}) \\
  &= 1.566\times(-0.2679)\times(-0.0543) \\
  &= 0.0228\;\text{rad} = \mathbf{1.3°} \quad\checkmark
\end{align*}}

\subsection{Solving for equilibrium speed $V$}

From Eq.\,\eqref{eq:fb2}:
\begin{equation}
  \boxed{
    V = \sqrt{\frac{m_0\,g}
    {\sin|\xi|\,(K_{D0}+K_D\alpha_d^2) + \cos|\xi|\,(K_{L0}+K_L\alpha_d)}}
  }
  \label{eq:Vequil}
\end{equation}

(Using $|\xi|$ and positive sign because we've already accounted for the direction
in the derivation.)

\insight{The speed equation reveals the two design levers:\\
\textbf{Numerator:} $m_0\,g$ is the net gravitational force.  More ballast $\rightarrow$ faster.
Speed scales as $\sqrt{m_0}$.\\
\textbf{Denominator:} the total hydrodynamic resistance.  Less drag $\rightarrow$ faster, but
also less lift $\rightarrow$ changes the equilibrium $\alpha$.}

\note{At $\xi=15°$, $\alpha_d=0.0228$\,rad, $m_0=0.120$\,kg:
\begin{align*}
  \text{drag term} &= \sin15°\times(1.083+25.06\times0.0228^2) = 0.2588\times1.096 = 0.284 \\
  \text{lift term} &= \cos15°\times(2.302+78.49\times0.0228) = 0.9659\times4.091 = 3.952 \\
  V &= \sqrt{0.120\times9.81/(0.284+3.952)} = \sqrt{1.177/4.236} = \sqrt{0.278} = \mathbf{0.527}\;\text{m/s} \quad\checkmark
\end{align*}}


%----------------------------------------------------------------------
\section{The complete equilibrium performance map}
\label{sec:perfmap}
%----------------------------------------------------------------------

Combining the equilibrium equations, we can tabulate performance across the full
operating envelope.  The horizontal speed is $V_h = V\cos\xi$.

\begin{table}[H]
\centering
\caption{Corrected steady glide equilibria for \seamoth{} v1 at $m_0=120$\,mL.}
\label{tab:equil}
\begin{tabular}{@{}r r r r r r r@{}}
\toprule
$\xi$ [°] & $\alpha_d$ [°] & $L/D$ & $V$ [m/s] & $V_h$ [m/s] & $V_z$ [m/s] & Re$_c$ \\
\midrule
10 & 3.1 & 5.67 & 0.42 & 0.41 & 0.07 & 33\,900 \\
12 & 2.2 & 4.70 & 0.47 & 0.46 & 0.10 & 37\,700 \\
15 & 1.3 & 3.73 & 0.53 & 0.51 & 0.14 & 42\,500 \\
18 & 0.8 & 3.08 & 0.58 & 0.55 & 0.18 & 46\,600 \\
20 & 0.5 & 2.75 & 0.61 & 0.57 & 0.21 & 49\,100 \\
25 & 0.0 & 2.14 & 0.68 & 0.61 & 0.29 & 54\,700 \\
\bottomrule
\end{tabular}
\end{table}

\begin{table}[H]
\centering
\caption{Speed at $\xi=15°$ for different ballast levels.}
\label{tab:speedballast}
\begin{tabular}{@{}r r r r@{}}
\toprule
$m_0$ [mL] & $m_0$ [kg] & $V$ [m/s] & $V_h$ [m/s] \\
\midrule
40  & 0.040 & 0.30 & 0.29 \\
80  & 0.080 & 0.43 & 0.42 \\
120 & 0.120 & 0.53 & 0.51 \\
160 & 0.160 & 0.61 & 0.59 \\
\bottomrule
\end{tabular}
\end{table}

\paragraph{Key observations.}
\begin{enumerate}[nosep]
  \item \textbf{Horizontal speed peaks at $\xi\approx20$--$25°$}, not $15°$ as
    previously stated.  The earlier spec had wrong coefficients that shifted the peak.
  \item \textbf{To reach $V_h>0.6$\,m/s} at $\xi=15°$, we need $m_0>180$\,mL---beyond
    our current buoyancy engine.  Options: steeper glide angle (20°, 0.57\,m/s at 120\,mL),
    or increase max ballast to 200+\,mL.
  \item \textbf{The SD7003 is never close to stall.}  Even at $\xi=10°$, $\alpha$
    is only 3.1°---well below the 8° stall angle.  This means we have substantial margin,
    and the MPC doesn't need an aggressive $\alpha$ limiter.
  \item \textbf{Re is above 30k for all operating points} except $m_0<50$\,mL at shallow
    angles.  The LSB risk is low in normal operation.
\end{enumerate}


%----------------------------------------------------------------------
\section{Stability of equilibrium glides}
\label{sec:stability}
%----------------------------------------------------------------------

Finding an equilibrium is only useful if the equilibrium is \emph{stable}---if the
glider naturally returns to the equilibrium after a disturbance.  Graver analyses
stability by linearizing the equations of motion around the equilibrium point
(G-\S4.1.3--4.1.5).

\subsection{Linearization}

Define small perturbations from the equilibrium: $\delta V = V - V_d$,
$\delta\xi = \xi - \xi_d$, $\delta\alpha = \alpha - \alpha_d$.  Graver shows that
the linearized longitudinal dynamics have the form:
\begin{equation}
  \dot{\vect{x}} = \mat{A}\,\vect{x} + \mat{B}\,\vect{u},
  \label{eq:linear}
\end{equation}
where $\vect{x}$ contains the perturbation states and $\mat{A}$ is the system matrix
whose eigenvalues determine stability.

\subsection{Static stability}

For \emph{static} stability (the initial tendency to return to equilibrium),
two conditions must be met:

\paragraph{Speed stability.}
If the glider is perturbed to a higher speed, drag should increase more than the
driving force decreases.  This is automatically satisfied when the drag polar has
$C_{D\alpha}>0$ (drag increases with $|\alpha|$), which is always true.

\paragraph{Pitch stability.}
If the glider pitches nose-up (increasing $\alpha$), a restoring nose-down moment
must develop.  This requires:
\begin{equation}
  C_{M\alpha} < 0
  \quad\text{and}\quad
  \frac{\partial M}{\partial\alpha} < 0,
  \label{eq:pitchstab}
\end{equation}
i.e., the aerodynamic centre must be \emph{behind} the centre of gravity.

\insight{This is the same stability criterion as for aircraft: aft CG $=$ stable.
For \seamoth{}, the wing is mounted at 55--65\% of body length from the nose.
Combined with the tail, this places the neutral point well behind the CG, giving
a positive \emph{static margin}.  The exact value depends on the hull's
destabilizing moment contribution and must be determined in CFD or tow-tank tests.}

\subsection{Dynamic stability}

The linearized system has characteristic modes analogous to aircraft:

\begin{description}[nosep]
  \item[Short period mode:] A fast, well-damped oscillation in $\alpha$ and pitch rate.
    This is stabilized by pitch damping ($C_{M_q}$, the pitch moment due to pitch rate)
    and is generally well-behaved for underwater gliders.
  \item[Phugoid mode:] A slow oscillation trading speed for altitude.  Graver analyses
    this in detail (G-\S4.2), adapting Lanchester's classical phugoid theory to underwater
    gliders.
\end{description}

\subsection{The phugoid mode for underwater gliders}
\label{sec:phugoid}

The phugoid is a long-period oscillation where the glider trades kinetic energy for
potential energy (like a roller-coaster).  For aircraft, the phugoid is lightly damped
and has period $T_p = \pi V_0\sqrt{2}/g$.

For underwater gliders, the phugoid is \emph{qualitatively different} from aircraft
(G-\S4.2).  The key differences:

\begin{enumerate}[nosep]
  \item \textbf{Buoyancy matters.}  In air, the aircraft weight is constant.  In water,
    the net driving force is $m_0\,g$, which is much smaller than $m_v\,g$ (since
    $m_0/m_v \approx 2\%$).  This makes the phugoid frequency much lower.
  \item \textbf{Added mass matters.}  The effective inertia of the glider includes
    the added mass of the surrounding water ($m_f$), increasing the effective mass
    by 10--30\%.  This further reduces the phugoid frequency.
  \item \textbf{$\alpha\neq0$ changes the dynamics.}  Graver shows (G-\S4.2.2) that
    nonzero equilibrium $\alpha$ breaks the symmetry of the phugoid phase portrait
    and can improve damping.
\end{enumerate}

\insight{For \seamoth{}, the phugoid is not a major concern.  Our dive depths are
only 5--30\,m, which is comparable to or shorter than one phugoid wavelength.  The
MPC will actively control through any phugoid oscillations rather than relying on
open-loop stability.  However, understanding the phugoid tells us the natural
timescale of speed/depth oscillations---this sets the minimum control bandwidth
the MPC must achieve.}


%----------------------------------------------------------------------
\section{Sawtooth gliding and inflections}
\label{sec:sawtooth}
%----------------------------------------------------------------------

Real glider operation involves repeated dive--climb cycles (``sawtooth'' pattern).
At the bottom and top of each sawtooth, the glider must \textbf{inflect}: reverse
its buoyancy and transition from downward to upward glide (or vice versa).

\subsection{The inflection problem}

During an inflection (G-\S4.1.6, \S6.1.5):
\begin{enumerate}[nosep]
  \item The pump reverses the ballast: $m_0$ goes from positive to negative.
  \item The pitch must reverse: from nose-down to nose-up.
  \item The speed may drop significantly if the transition is slow.
  \item The glider briefly passes through zero speed---where hydrodynamic control
    authority vanishes.
\end{enumerate}

The inflection is the most dynamic and challenging part of the flight.  Graver
simulates this and shows that with proper control, the transition can be smooth
(G-Fig.\,4.6--4.8).

\subsection{Inflection time for \seamoth{}}

The inflection time is set by two factors:

\paragraph{Pump time.}
To reverse from $m_0=+120$\,mL to $m_0=-120$\,mL requires pumping 240\,mL of oil.
At 500\,mL/min loaded flow: $t_\text{pump} = 240/500\times60 = 28.8$\,s.

\paragraph{Pitch transition.}
With VBC, the pitch moment depends on the pressure differential between forward and aft
bladders.  The moment arm is roughly half the hull length ($\sim$0.35\,m), and the
restoring force is the buoyancy differential.  A rough estimate: with 40\,mL differential
(20\,mL extra forward, 20\,mL less aft), the restoring force is $0.040\times9.81\approx0.4$\,N
at 0.35\,m arm $= 0.14\;\text{N}\cdot\text{m}$ of pitch moment.

Against a pitch inertia of $I_\theta\approx0.3\;\text{kg}\cdot\text{m}^2$ (estimated),
the angular acceleration is $\sim$0.47\,rad/s$^2$ $\approx$ 27°/s$^2$.  A 30° pitch
reversal thus takes $\sim$1.5\,s.

The total inflection time is dominated by the pump: $\sim$30\,s.  This is consistent
with the 3--7\,s estimate in the spec only if a dual-pump configuration is used.
With a single pump, expect 25--35\,s inflections.


%======================================================================
%  PART IV: THREE-DIMENSIONAL MOTION
%======================================================================
\part{Three-Dimensional Motion}

%----------------------------------------------------------------------
\section{Turning}
\label{sec:turning}
%----------------------------------------------------------------------

An underwater glider turns by \emph{banking}: rolling to one side so that the
lift vector has a horizontal component.  This is identical to how aircraft turn.

\subsection{The turning mechanism}

In conventional gliders (Slocum, Seaglider), roll is induced by laterally shifting
the internal sliding mass or by using a rudder.  In \seamoth{}, roll is induced by
differential bladder inflation (port vs.\ starboard), which shifts the centre of
buoyancy laterally.

Once banked at roll angle $\phi$, the horizontal component of lift provides the
centripetal force for turning:
\begin{equation}
  L\sin\phi = \frac{m_v\,V^2}{R},
  \label{eq:turn}
\end{equation}
where $R$ is the turn radius.  Solving:
\begin{equation}
  R = \frac{m_v\,V^2}{L\sin\phi}.
  \label{eq:turnrad}
\end{equation}

\note{At $V=0.53$\,m/s, $\phi=30°$, $L\approx m_0 g/\cos\xi = 0.120\times9.81/\cos15° = 1.22$\,N:
\[R = \frac{8.5\times0.53^2}{1.22\times\sin30°} = \frac{2.39}{0.61} = 3.9\;\text{m}.\]
This meets the $<$5\,m turning radius target.  More aggressive banking (45°) gives $R\approx2.8$\,m.}

\subsection{Roll--yaw coupling}
\label{sec:rollyaw}

An important subtlety (G-\S4.3.5): rolling the glider also creates a yaw moment.
This happens because rolling changes the angle of attack differently on the two wings
(the ``inside'' wing sees higher effective $\alpha$ than the ``outside'' wing during
the roll maneuver), creating differential drag and thus a yaw moment.

For \seamoth{}, the wing position relative to the CG determines whether this coupling
is \emph{proverse} (yaw in the desired turn direction) or \emph{adverse} (yaw opposing
the turn).  Wings behind the CG give proverse coupling; wings ahead give adverse.
Our wing is at 55--65\% from the nose---likely slightly aft of CG, which is favorable.

\subsection{Spiral glides}

A glider can fly a steady turning descent (or ascent) called a \textbf{spiral glide}
(G-\S4.3.2).  This is a useful survey pattern: the glider descends in a helix,
imaging the reef from multiple angles.

The equilibrium spiral combines straight-glide equations with the turning
equation.  Graver derives the full 3D equilibrium conditions (G-\S4.3.1--4.3.2)
and shows that spiral glides exist over a range of bank angles and glide angles.

\note{For reef survey, a descending spiral at $\xi=15°$, $\phi=20°$, $V=0.53$\,m/s
gives: turn radius $\approx5.8$\,m, vertical descent rate $\approx0.14$\,m/s,
circumference $\approx36$\,m, time per revolution $\approx68$\,s.  In a 15\,m dive,
the glider completes $\sim$1.6 revolutions---enough for useful photogrammetric coverage
of a coral bommie.}


%======================================================================
%  PART V: DESIGN EQUATIONS
%======================================================================
\part{Design Equations}

%----------------------------------------------------------------------
\section{Speed, ballast, and energy}
\label{sec:design}
%----------------------------------------------------------------------

Graver Chapter~7 applies the equilibrium equations to analyse how design parameters
affect glider performance.  This is the most directly useful section for \seamoth{}
sizing.

\subsection{Speed vs.\ ballast: the $\sqrt{m_0}$ law}

From Eq.\,\eqref{eq:Vequil}, at fixed $\xi$ and $\alpha_d$:
\begin{equation}
  V \propto \sqrt{m_0}.
  \label{eq:Vsqrt}
\end{equation}

Doubling the ballast increases speed by $\sqrt{2}\approx41\%$.  This is confirmed
by Table~\ref{tab:speedballast}: $V(160)/V(80) = 0.61/0.43 = 1.42$.

\subsection{Speed vs.\ vehicle size: the $\sqrt{m_v\beta}$ law}

For a family of geometrically similar vehicles with ballast fraction $\beta=m_0/m_v$
(G-\S7.2.3):
\begin{equation}
  V \propto \sqrt{m_v\,\beta}.
  \label{eq:Vscale}
\end{equation}

A larger vehicle with the same $\beta$ goes faster simply because it has more ballast.
This is why the 990\,kg STERNE glider achieves 1.3\,m/s while the 52\,kg Slocum achieves
0.4\,m/s, despite similar $\beta$.

\subsection{Ballast fraction comparison}

\begin{table}[H]
\centering
\caption{Ballast fraction and speed for existing gliders.}
\begin{tabular}{@{}l c c c c@{}}
\toprule
Vehicle & $m_v$ [kg] & $\beta$ [\%] & $V$ [m/s] & Mission \\
\midrule
Seaglider & 52 & 0.8 & 0.25 & Trans-ocean \\
Slocum Electric & 52 & 1.5 & 0.40 & Coastal \\
STERNE & 990 & 1.2 & 1.30 & Near-seabed survey \\
\seamoth{} v1 & 8.5 & 1.88 & 0.5--0.6 & Reef survey \\
\bottomrule
\end{tabular}
\end{table}

\seamoth{}'s higher $\beta$ compensates for its small size.  We intentionally
sacrifice endurance (more pumping energy per unit distance) for speed, because
reef survey requires covering area quickly during daylight hours.

\subsection{Optimal glide angle for maximum horizontal speed}

The horizontal speed is $V_h = V\cos\xi$.  From Eq.\,\eqref{eq:Vequil},
$V$ increases with $|\xi|$ (steeper $=$ faster) but $\cos\xi$ decreases.  The
product $V_h$ has a maximum at some optimal $\xi^*$.

Graver shows (G-\S7.2.2) that this optimal angle depends on the drag polar but
is typically in the range 15--25° for conventional gliders.  For \seamoth{}, the
corrected equilibrium table shows $V_h$ peaks around $\xi\approx20$--$25°$.

\insight{The original spec targeted $\xi=15°$ based on incorrect coefficients.
With the corrected hull drag, the optimal angle shifts to $\sim$20°.  This is
actually \emph{good news} for reef survey: steeper angles mean more vertical
coverage per horizontal distance, which improves photogrammetric overlap.  The
trade-off is that you reach the bottom sooner and need more inflections per
kilometer of horizontal distance.}


%----------------------------------------------------------------------
\section{Energy and endurance}
\label{sec:energy}
%----------------------------------------------------------------------

\subsection{Pump energy per cycle}

Each dive--climb cycle requires reversing the buoyancy.  The pump moves a volume
$\Delta\mathcal{V} = 2m_0/\rho$ of oil against the hydrostatic pressure.
Over a sawtooth dive to depth $d$, the average pressure is:
\begin{equation}
  \bar{P} = \frac{\rho\,g\,d}{2}.
  \label{eq:Pavg}
\end{equation}

The mechanical work per cycle is:
\begin{equation}
  E_\text{mech} = \Delta\mathcal{V}\times\bar{P}
  = \frac{2m_0}{\rho}\times\frac{\rho\,g\,d}{2}
  = m_0\,g\,d.
  \label{eq:Emech}
\end{equation}

\insight{The $\rho$ cancels!  The pump energy is exactly the potential energy of
lifting mass $m_0$ through height $d$.  This is physically obvious in retrospect:
the glider converts gravitational potential energy into kinetic energy of gliding;
to reverse the cycle, you must put that energy back.}

The electrical energy, accounting for pump efficiency:
\begin{equation}
  E_\text{elec} = \frac{E_\text{mech}}{\eta_\text{pump}}.
  \label{eq:Eelec}
\end{equation}

\note{$m_0=0.120$\,kg, $d=15$\,m (average sawtooth depth to 30\,m),
$\eta=0.45$:
\begin{align*}
  E_\text{mech} &= 0.120\times9.81\times15 = 17.7\;\text{J}, \\
  E_\text{elec} &= 17.7/0.45 = 39.2\;\text{J} = 0.0109\;\text{Wh per cycle}.
\end{align*}}

\subsection{System power budget}

The total average power consumption:
\begin{equation}
  P_\text{total} = P_\text{hotel} + P_\text{pump,avg} + P_\text{SBC,avg},
  \label{eq:Ptotal}
\end{equation}
where:
\begin{itemize}[nosep]
  \item $P_\text{hotel}$: always-on loads (H755 + sensors) $\approx 0.8$\,W.
  \item $P_\text{pump,avg} = E_\text{elec}\times f_\text{cycles}$: pump energy $\times$
    cycling frequency.  At 6 cycles/hr: $P_\text{pump}=39.2\times6/3600=0.065$\,W.
  \item $P_\text{SBC,avg} = P_\text{active}\times\delta + P_\text{sleep}\times(1-\delta)$:
    Jetson duty-cycled.  At $\delta=8\%$: $P_\text{SBC}=7.0\times0.08+0.3\times0.92=0.84$\,W.
\end{itemize}

\note{$P_\text{total} = 0.8 + 0.065 + 0.84 \approx 1.7$\,W.  This is lower than the
2.4\,W in the spec because the corrected coefficients give lower speeds (less pump work)
and we're using a more careful power breakdown.  The spec's 2.4\,W included valve power
and DC-DC losses that should still be added for the full system.  With those:
$P_\text{total}\approx2.2$\,W.}

\subsection{Endurance and range}

\begin{equation}
  T = \frac{E_\text{battery}}{P_\text{total}}, \qquad
  R = V_h\times T.
  \label{eq:endurance}
\end{equation}

\note{$E_\text{batt}=116$\,Wh, $P_\text{total}=2.2$\,W, $V_h=0.51$\,m/s:
\begin{align*}
  T &= 116/2.2 = 52.7\;\text{hours} \approx 2.2\;\text{days}, \\
  R &= 0.51\times52.7\times3.6 = 97\;\text{km}.
\end{align*}
Active survey at 30\% Jetson duty: $P\approx3.5$\,W, $T\approx33$\,h, $R\approx61$\,km.}


%----------------------------------------------------------------------
\section{Wing sizing considerations}
\label{sec:wingsizing}
%----------------------------------------------------------------------

Graver analyses the effect of wing area on L/D (G-\S7.3, Fig.\,7.5).  The key
trade-off:

\begin{itemize}[nosep]
  \item \textbf{Larger wing:} more lift per unit $\alpha$ $\rightarrow$ lower $\alpha$
    at a given $\xi$ $\rightarrow$ lower induced drag.  But also more wetted area
    $\rightarrow$ more profile drag.
  \item \textbf{Smaller wing:} less profile drag but higher $\alpha$ $\rightarrow$
    higher induced drag and potentially stall at steep angles.
\end{itemize}

For \seamoth{}, since induced drag is only 1\% of total, \emph{reducing wing area
would actually reduce total drag} (less profile drag) with negligible increase in
induced drag.  However, smaller wings reduce roll authority in the VBC scheme
(less moment arm for lateral CB shift) and reduce the wing's contribution to pitch
stability.

Our 300\,cm$^2$ wing area is a reasonable compromise.  Sensitivity analysis in the
notebook shows that L/D varies only 5--10\% over the range 200--400\,cm$^2$.


%======================================================================
%  PART VI: PARAMETER IDENTIFICATION
%======================================================================
\part{Parameter Identification}

%----------------------------------------------------------------------
\section{What to measure in pool tests}
\label{sec:paramid}
%----------------------------------------------------------------------

Graver Chapter~5 describes how to identify hydrodynamic coefficients from flight
data.  The basic method:

\subsection{Steady glide method}

\begin{enumerate}[nosep]
  \item Fly the glider at several known $m_0$ values.
  \item At each $m_0$, let the glider reach steady state and measure:
    \begin{itemize}[nosep]
      \item Pitch angle $\theta$ (from IMU).
      \item Depth rate $\dot{z}$ (from depth sensor differentiation).
      \item Speed $V$ (from depth rate and pitch: $V=\dot{z}/\sin\xi$,
        or better, from a DVL if available).
    \end{itemize}
  \item Compute $\xi = \arctan(\dot{z}/\dot{x})$ and $\alpha = \theta - \xi$.
  \item From Eq.\,\eqref{eq:DoverL}: $D/L = -\tan\xi$ gives the drag-to-lift ratio
    at each $\alpha$.
  \item From Eq.\,\eqref{eq:Vequil}: knowing $V$, $\xi$, and $m_0$, solve for
    the $K$ parameters.
  \item Fit the coefficient model ($C_{D0}$, $C_{D\alpha}$, $C_{L0}$, $C_{L\alpha}$)
    to the data.
\end{enumerate}

\warning{The main difficulty is measuring speed $V$ accurately.  Without a DVL,
$V$ must be inferred from $\dot{z}$ and $\theta$, which requires accurate pitch
measurement and the assumption of steady state.  In pool tests with GPS-denied
conditions, video tracking from above may be the most practical speed measurement.}

\subsection{What pool tests will tell us}

\begin{enumerate}[nosep]
  \item \textbf{True $C_{D0}$:} Our estimate of 0.071 is based on component buildup.
    Real vehicles always have higher drag than predicted (cables, connectors, imperfect
    fairing, hull roughness).  Expect 10--30\% higher: $C_{D0}\approx0.08$--$0.09$.
  \item \textbf{True neutral buoyancy trim:} Even small trim errors ($\pm$5\,g) affect
    the equilibrium.  The pool test calibrates the relationship between pump command and
    actual $m_0$.
  \item \textbf{Inflection dynamics:} How fast can we actually reverse pitch?  What is
    the speed loss during inflection?
  \item \textbf{VBC effectiveness:} How much differential bladder inflation is needed for
    a given pitch/roll rate?
  \item \textbf{Turn performance:} Actual turn radius vs.\ bank angle.
\end{enumerate}


%----------------------------------------------------------------------
\section{How to read Graver's dissertation}
\label{sec:readingguide}
%----------------------------------------------------------------------

If you want to go deeper into specific topics beyond this document, here is a
prioritized reading guide:

\begin{table}[H]
\centering
\caption{Reading guide for Graver (2005), prioritized for \seamoth{}.}
\small
\begin{tabular}{@{}clp{6.5cm}c@{}}
\toprule
Priority & Graver section & What it covers & Pages \\
\midrule
1 & \S4.1 (pp.\,86--110) & Longitudinal equilibria, stability, sawtooth sim & 24 \\
2 & \S7.1--7.3 (pp.\,224--243) & Gliders vs sailplanes, speed/sizing equations & 19 \\
3 & \S3.2.7 (pp.\,82--85) & Hydrodynamic force model & 4 \\
4 & \S5.1--5.2 (pp.\,147--177) & Slocum model and parameter ID from sea data & 30 \\
5 & \S4.3 (pp.\,125--144) & 3D gliding, spirals, turning, roll-yaw & 19 \\
6 & \S4.2 (pp.\,110--121) & Phugoid mode analysis & 11 \\
7 & \S6.1 (pp.\,181--196) & Controller design (linear observer/controller) & 15 \\
\midrule
& & \textbf{Total for priorities 1--5} & \textbf{96 pages} \\
\bottomrule
\end{tabular}
\end{table}

\textbf{Skip entirely:} Chapter~1 (intro), Chapter~2 (history---interesting but
not technical), \S3.1 (kinematics---covered here), \S3.2.1--3.2.6 (detailed EOM
derivation---important but you can trust the results), \S6.2--6.4 (Slocum-specific
control), \S7.4 (flying wing designs), Chapter~8 (conclusions), Appendix.


%======================================================================
\section*{Summary}
%======================================================================

The five equations that govern \seamoth{}'s steady-state flight:

\begin{enumerate}
  \item \textbf{Glide ratio} (Eq.\,\ref{eq:DoverL}):
    $D/L = -\tan\xi$.

  \item \textbf{Equilibrium $\alpha$} (Eq.\,\ref{eq:alpha_d}):
    quadratic in $\alpha$ from the force balance; gives $\alpha_d=1.3°$ at $\xi=15°$.

  \item \textbf{Equilibrium speed} (Eq.\,\ref{eq:Vequil}):
    $V=\sqrt{m_0 g/\text{(hydro resistance)}}$; gives $V=0.53$\,m/s at 120\,mL.

  \item \textbf{Pump energy} (Eq.\,\ref{eq:Emech}):
    $E=m_0 g d$; gives 17.7\,J mechanical per cycle.

  \item \textbf{Endurance} (Eq.\,\ref{eq:endurance}):
    $T=E_\text{batt}/P_\text{total}$; gives $\sim$53\,h continuous.
\end{enumerate}

Everything in the Jupyter notebook is a numerical evaluation of these five equations
with \seamoth{} parameters.

\vspace{1em}\hrule\vspace{1em}

\section*{References}
\begin{enumerate}[label={[\arabic*]},nosep]
  \item J.\,G.\ Graver, \emph{Underwater Gliders: Dynamics, Control and Design},
    Ph.D.\ dissertation, Princeton University, 2005.
    \url{https://naomi.princeton.edu/wp-content/uploads/sites/744/2021/03/jggraver-thesis-4-11-05.pdf}
  \item S.\,Selig, J.\,Donovan, D.\,Fraser, \emph{Airfoils at Low Speeds}, Soartech~8, 1989.
  \item K.\,Galloway et al., ``ReefGlider: A Vectored Buoyancy Controlled Underwater
    Glider for Reef Survey,'' arXiv:2405.06033, 2024.
  \item T.\,Fossen, \emph{Handbook of Marine Craft Hydrodynamics and Motion Control},
    Wiley, 2011.
  \item M.\,S.\ Triantafyllou and F.\,S.\ Hover, ``Maneuvering and Control of Marine
    Vehicles,'' MIT OCW 2.154, 2004.
\end{enumerate}

\end{document}
